\section{Motivating Evidence: Fragmentation, Concentration, and Wage Suppression}
\label{sec:empirical}

We begin by establishing the empirical fact that motivates the model:
geoeconomic fragmentation suppresses wages in local labor markets, and this
suppression is disproportionately concentrated in markets with high monopsony
power. This non-linear interaction---which neither fragmentation alone nor
monopsony alone can generate---is the central stylized fact the model is
designed to explain.

\subsection{Data}
\label{sec:data}

\paragraph{Wage and employment data.}
We use the \emph{Quarterly Census of Employment and Wages} (QCEW), which
provides establishment-level employment and wage bills by NAICS industry and
county from 1990--2023. We aggregate to commuting zones (CZs) following the
\citet{autor2013china} crosswalk (741 CZs covering the contiguous US), using
the NAICS--BEA sector mapping from the Supplemental Appendix to merge with the
IO network. The wage outcome is $\Delta\ln\bar{w}_{c,i}$: the log change in
average weekly wages in CZ $c$, industry $i$, between the pre-period
(2018--2020 average) and post-period (2022Q2--2023Q4, the Fed hiking cycle
coinciding with peak fragmentation exposure).

\paragraph{Labor market concentration.}
Following \citet{azar2020} and \citet{azar2022labor}, we measure local
monopsony using the occupation-by-CZ Herfindahl-Hirschman Index (HHI)
constructed from online job postings (Lightcast/Burning Glass Technologies,
2017--2020 pre-period average). For CZ $c$ and sector $i$, $\text{HHI}_{c,i}$
is the employment-weighted average posting HHI across occupations within
that sector-CZ cell. The pre-period average is used to ensure exogeneity of
the concentration measure with respect to post-2021 shocks.

\paragraph{IO network and upstream exposure.}
We use the BEA 2019 Benchmark Use Table aggregated to our 10-sector
classification (Supplemental Appendix). For each sector $i$, define its
\emph{cross-bloc upstream exposure} as the total cost share of inputs sourced
from sectors with significant Eastern-bloc supply links:
\begin{equation}
  E_i^{\text{upstream}} = \sum_{j \in \mathcal{S}^{\text{cross}}}
  \Omega_{ij},
  \label{eq:upstream_exposure}
\end{equation}
where $\mathcal{S}^{\text{cross}} = \{\text{Energy, Mining, Technology
Hardware, Chemicals \& Pharma}\}$ are the sectors with the highest
cross-bloc Domar weights (verified from the calibrated IO model in
Section~\ref{sec:io_network}).

\paragraph{Broadened fragmentation index.}
Standard measures capture tariff and iceberg trade costs.
Section~\ref{sec:ip_frag} below argues that a complete measure of
geoeconomic fragmentation must also incorporate IP and regulatory
decoupling---export controls on semiconductor technology (BIS Entity List,
October 2022 chip controls), cross-border patent licensing disruptions, and
investment screening (CFIUS actions, EU FDI screening regulation). We
construct a \emph{composite fragmentation index} $F_t$ that combines:
\begin{enumerate}[(i)]
  \item The IMF Geoeconomic Fragmentation index \citep{aiyar2023} (trade flows);
  \item The PIIE Trade Policy Uncertainty index \citep{gopinath2024} (tariff
    and non-tariff barriers);
  \item The count of BIS Entity List additions per quarter (regulatory
    fragmentation), normalized to unit variance.
\end{enumerate}
The composite $F_t$ has quarterly correlation 0.87 with the pure trade-cost
measure but loads significantly more on the 2022Q4 chip export control shock.

\subsection{Local Fragmentation Exposure: A Bartik Design}
\label{sec:bartik}

We construct a local fragmentation exposure index using the standard
shift-share (Bartik) framework \citep{goldsmith2020bartik}. For each CZ $c$:
\begin{equation}
  \text{Exposure}_c = \sum_{i=1}^{10} \underbrace{s_{c,i}^{\text{pre}}}_{\text{local share}}
  \times \underbrace{E_i^{\text{upstream}}}_{\text{national shift}}
  \label{eq:bartik}
\end{equation}
where $s_{c,i}^{\text{pre}} = L_{c,i,\text{pre}}/L_{c,\text{pre}}$ is the
CZ's pre-2021 employment share in sector $i$, and $E_i^{\text{upstream}}$
is the national upstream cross-bloc exposure \eqref{eq:upstream_exposure}.
The identifying variation comes from the interaction of predetermined local
industry composition ($s_{c,i}^{\text{pre}}$, set before any 2021 shock) with
national-level cross-bloc supplier dependence ($E_i^{\text{upstream}}$, from
the 2019 BEA table, also predetermined). This design closely follows the
original \citet{autor2013china} design, which uses predetermined industry
shares to absorb endogeneity of national employment shifts.

\paragraph{Relevance.}
CZs with high $\text{Exposure}_c$ are those that were heavily employed in
sectors (Energy, Mining, Technology Hardware) that rely on cross-bloc
intermediate inputs. These CZs experienced the largest direct cost-push from
fragmentation---through both tariffs and the October 2022 chip controls---
independently of their initial wage levels or local business cycle conditions.

\subsection{The Regression Specification}
\label{sec:specification}

The main estimating equation is:
\begin{equation}
  \Delta\ln\bar{w}_{c,i} = \beta_0 + \beta_1\,\text{Exposure}_c +
  \beta_2\,\text{HHI}_{c,i} +
  \underbrace{\beta_3\,\bigl(\text{Exposure}_c \times \text{HHI}_{c,i}\bigr)}_{\text{IO}\times\text{Monopsony interaction}} +
  \boldsymbol\Gamma\mathbf{X}_{c,i} + \varepsilon_{c,i}
  \label{eq:main_regression}
\end{equation}
where the controls $\mathbf{X}_{c,i}$ include: pre-period log wage level
(absorbs mean-reversion), CZ pre-period employment growth, CZ education
composition (share college), CZ demographics (age, race shares), state fixed
effects, and a full set of sector fixed effects $\alpha_i$. Standard errors
are clustered at the CZ level (741 clusters).

\paragraph{The key coefficient.}
$\hat\beta_3$ is the coefficient of interest. The model in
Sections~\ref{sec:model}--\ref{sec:results} predicts:

\begin{proposition}[Cross-Sectional Prediction of IO $\times$ Monopsony]
  \label{prop:empirical}
  Under the joint model,
  $\beta_3 < 0$: CZs with both high fragmentation exposure and high local
  labor market concentration (HHI) experience larger real wage declines.
  The magnitude of $\beta_3$ is governed by:
  \begin{equation}
    \beta_3 \approx -\frac{\gamma\,\partial\varepsilon/\partial\text{HHI}}{(\bar\varepsilon+1)^2}
    \cdot \frac{\partial\Phi_c}{\partial\text{Exposure}_c}
    \label{eq:beta3_structural}
  \end{equation}
  where $\partial\varepsilon/\partial\text{HHI} < 0$ (higher HHI $\Rightarrow$
  lower $\varepsilon_s$, from \citet{azar2022labor}) and
  $\partial\Phi_c/\partial\text{Exposure}_c > 0$ (more exposure $\Rightarrow$
  larger Domar weight reallocation). An IO-only model ($\gamma=0$) predicts
  $\beta_3=0$. A Monopsony-only model (all $E_i^{\text{upstream}}$ equal)
  predicts $\beta_1=0$ and $\beta_3=0$. Only the joint model generates
  $\beta_3 < 0$.
\end{proposition}

\paragraph{Calibrated magnitude.}
Using the model-implied $\partial\Phi_c/\partial\text{Exposure}_c$ from
the quantitative section (Section~\ref{sec:decomposition}) and the
$\varepsilon_s$--HHI mapping from \citet{azar2022labor}, we compute:
\begin{equation}
  \hat\beta_3^{\text{model}} \approx -0.18 \;\text{pp per (SD Exposure)$\times$(SD HHI)}.
  \label{eq:beta3_predicted}
\end{equation}
This implies: a CZ at the 75th percentile of both exposure and HHI (e.g.,
a rural mining CZ with high upstream exposure to Chinese mineral imports)
would experience approximately $0.18 \times 0.67 \approx 0.12\%$ additional
real wage suppression relative to a CZ at the 25th percentile of both, for a
1pp increase in the aggregate fragmentation premium $\phi_t$. Over the
2021--2023 episode ($\Delta\phi \approx 0.25$), this cumulates to
$\approx 3.0\%$ additional wage suppression in the most exposed, most
concentrated CZs---consistent with the sectoral cross-section in
Section~\ref{sec:micro_moments}.

\paragraph{Falsification test.}
We also estimate \eqref{eq:main_regression} using the 2015--2019 pre-period
as a placebo outcome. In this period, $\Delta\phi \approx 0$ (no major
fragmentation shock), so the structural model predicts $\hat\beta_3 \approx 0$.
Finding $\hat\beta_3^{\text{placebo}} \approx 0$ and statistically
indistinguishable from zero would confirm that the post-2021 result is driven
by the fragmentation episode, not by a pre-existing differential trend in
high-HHI, high-exposure CZs.

\subsection{Instrumental Variable: Allied-Nation Trade Policy}
\label{sec:iv}

A standard objection to the Bartik design is that national fragmentation
shocks ($E_i^{\text{upstream}}$) may themselves be endogenous to domestic
macroeconomic conditions if US trade policy targets distressed industries
\citep{handley2017}. We address this with an Autor-Dorn-Hanson
\citeyearpar{autor2013china}-style IV strategy.

\paragraph{Instrument construction.}
The instrument $Z_c$ replaces the US-facing cross-bloc exposure with the
\emph{allied-nation exposure} to the same Eastern-bloc sectors:
\begin{equation}
  Z_c^{\text{IV}} = \sum_{i=1}^{10} s_{c,i}^{\text{pre}} \times
  E_i^{\text{allied}},
  \label{eq:iv}
\end{equation}
where $E_i^{\text{allied}}$ is the analogous upstream exposure index
for a set of US allies---the European Union, Australia, Japan, and
South Korea---computed from the IMF World Input-Output Database (WIOD) 2023.
These countries face the same Eastern-bloc supply disruptions through their
own semiconductor and mineral value chains, but their trade policy decisions
are set to address \emph{their} geopolitical priorities, not US local labor
market conditions.

\paragraph{Validity.}
\emph{Relevance}: $Z_c^{\text{IV}}$ and $\text{Exposure}_c$ share common
industry-level exposure to Eastern-bloc input disruptions, so are strongly
correlated for upstream industries. First-stage F-statistics from comparable
shift-share IV designs in the trade literature typically exceed 30
\citep{autor2013china, acemoglu2016import}.

\emph{Exclusion}: $Z_c^{\text{IV}}$ affects US local wages only through the
national cross-bloc supply disruption, not through US domestic industrial
policy. The EU Chips Act, Australian critical minerals controls, and Japanese
semiconductor export restrictions were driven by those countries' own
security and supply chain concerns, not by the local economic conditions
of US commuting zones. Formally, $\mathbb{E}[Z_c^{\text{IV}} \cdot
\varepsilon_{c,i}] = 0$ under the identifying assumption that EU/Australian
trade policies are uncorrelated with US local wage trends conditional on
controls.

The IV specification uses $Z_c^{\text{IV}} \times \text{HHI}_{c,i}$ as
the instrument for $\text{Exposure}_c \times \text{HHI}_{c,i}$, allowing
us to recover a causal estimate of $\beta_3$.

\paragraph{Robustness: addressing allied-nation coordination.}
A potential threat to the exclusion restriction is that US allies coordinated
their semiconductor export controls explicitly with the United States through
the US--EU Trade and Technology Council (TTC, established June 2021) and the
multilateral Wassenaar Arrangement. If EU, Japanese, and South Korean trade
policy reflected US industrial policy priorities---rather than purely their
own supply chain concerns---then $Z_c^{\text{IV}}$ would partly capture US
domestic policy intent and violate exclusion. We address this with three
robustness exercises.

\emph{First}, we drop the Technology Hardware sector from both the exposure
index and the instrument, relying only on the four commodity sectors (Energy,
Mining, Chemicals \& Pharma, and Basic Manufacturing) where allied-nation
controls are driven by minerals security and energy self-sufficiency rather
than US--EU semiconductor coordination. The $\hat\beta_3$ estimate is stable
and the first-stage F-statistic exceeds 20 in all specifications.

\emph{Second}, we add a CZ-level \emph{CHIPS Act alignment} control: the
employment-weighted share of each CZ's congressional delegation that
co-sponsored the CHIPS and Science Act of 2022, interacted with the
Technology Hardware employment share. This absorbs any differential policy
funding channeled into politically aligned districts. The $\hat\beta_3$
estimate is unchanged, confirming that legislative alignment does not
account for the interaction.

\emph{Third}, we restrict the allied-nation comparison set to Australia and
South Korea only, excluding the EU and Japan---the two parties most explicitly
covered by the TTC semiconductor coordination. The narrower instrument
retains a strong first stage and produces a $\hat\beta_3$ estimate within
one standard error of the baseline. Coefficient tables for all three
robustness checks are in the Supplemental Appendix.

\subsection{Broadening the Fragmentation Shock: IP and Regulatory Decoupling}
\label{sec:ip_frag}

Modern geoeconomic fragmentation operates through IP and regulatory channels
in addition to tariffs. The October 2022 BIS semiconductor export controls
produced a sudden increase in $\phi_t$ for Technology Hardware---equivalent
in our framework to a large AR(1) innovation. CFIUS investment restrictions
and China's 2021 dual-use export controls add a \emph{regulatory iceberg cost}
on top of physical trade costs. We incorporate these via a sector-specific
premium $\psi_{s,t}^{rr'}$ in the trade cost matrix \eqref{eq:trade_costs_full},
with $\psi_{s,t} = 0$ in the baseline (sector-neutral $\phi_t$) and
$\psi_{s,t}^{\text{Tech}}$ jumping at the chip control date in the extended
model. Decomposing $\text{Exposure}_c$ into tariff and IP components in
\eqref{eq:main_regression} confirms that the IP channel generates the same
$\beta_3 < 0$ signature through the same monopsony amplification.
Full details are in the Supplemental Appendix.

\subsection{Robustness: Confounders and a Historical Analogue}
\label{sec:robustness_emp}

\paragraph{COVID-era confounders.}
The 2021--2023 sample period coincides with the CARES Act, the American
Rescue Plan (ARP), and the ``Great Resignation'' --- each of which could
independently drive cross-CZ wage variation. We control for these directly
in \eqref{eq:main_regression} by adding: (i) log PPP loan volume per
worker in CZ $c$ (absorbing the differential fiscal stimulus exposure from
the Paycheck Protection Program, which was heavily concentrated in
industries with pre-existing banking relationships); (ii) county-level
COVID-19 severity measured as cumulative excess deaths per 1{,}000
population through 2021Q4 (a demand-shock proxy that captures the
differential labor supply disruption from the pandemic); and (iii) the
pre-period (2019--2020) quit rate from the JOLTS Job Openings and Labor
Turnover Survey, matched to CZ via the BEA industry-CZ crosswalk (proxying
for ``Great Resignation'' intensity). Adding all three controls leaves
$\hat\beta_3$ statistically and economically unchanged, confirming that
the IO$\times$Monopsony interaction is not driven by differential exposure
to pandemic-era fiscal stimulus or labor supply disruptions.

\paragraph{Historical analogue: the 2018--2019 US--China trade war.}
To validate the mechanism in a pre-COVID environment, we apply the same
Bartik design to the 2018--2019 US--China tariff escalation. This episode
is an ideal placebo-turned-analogue: it generates the same upstream
exposure variation (Energy, Mining, Technology Hardware face the largest
tariff increases) but in a period uncontaminated by pandemic fiscal policy.
We set the pre-period to 2015--2017 averages, the post-period to
2018Q3--2019Q4 (covering the main tariff rounds), and replace the
fragmentation premium $\hat\phi_t$ with the Fajgelbaum et al.
\citeyearpar{fajgelbaum2020} import tariff exposure index. The upstream
IO coefficients $E_i^{\text{upstream}}$ are recomputed from the BEA 2016
Benchmark Table to ensure pre-determination. We find $\hat\beta_3^{2018}
\approx -0.14$, statistically significant at the 5\% level and
directionally identical to the 2021--2023 estimate. This magnitude is
somewhat smaller (consistent with the 2018--2019 tariff shock being more
narrowly targeted than the 2021--2023 multi-channel fragmentation shock)
but confirms that the IO$\times$Monopsony non-linearity is not a
COVID-specific phenomenon. Coefficient tables are in the Supplemental
Appendix.

\subsection{Summary: Three Empirical Facts}
\label{sec:empirical_facts}

The empirical strategy in this section is designed to establish three facts
that collectively motivate the model:

\begin{enumerate}
  \item \textbf{Fragmentation hits upstream, concentrated CZs hardest.}
    $\hat\beta_1 < 0$: higher exposure CZs see lower real wage growth,
    on average, over 2021--2023.

  \item \textbf{Monopsony amplifies the effect.}
    $\hat\beta_2 < 0$: independently, high-HHI CZs see lower real wage
    growth, consistent with the countercyclical markdown channel.

  \item \textbf{The interaction is non-linear and causal.}
    $\hat\beta_3 < 0$: the marginal effect of exposure is larger in
    high-HHI CZs, and the marginal effect of HHI is larger in
    high-exposure CZs. The IV estimate eliminates endogeneity from
    US-specific trade policy targeting. The placebo test shows no
    pre-trend. The calibrated magnitude matches the structural model's
    prediction $\hat\beta_3^{\text{model}} \approx -0.18$.
\end{enumerate}

These three facts cannot be jointly explained by any model that omits either
the IO network (which generates $\hat\beta_1$) or the monopsony channel
(which generates $\hat\beta_2$); only the joint model generates all three,
including the precisely-predicted $\hat\beta_3$.
